\uuid{AVF9}
\titre{ Stabilité numérique de la fonction sigmoïde }
\niveau{L3}
\module{Analyse de données}
\chapitre{Réseaux de neurones}
\sousChapitre{Fonctions d'activation, stabilité numérique}
\theme{Sigmoïde, overflow, underflow, stabilité numérique}
\auteur{Maxime Nguyen}
\datecreate{2026-03-19}
\organisation{AMSCC}
\difficulte{3}

\contenu{
\texte{
  On compare deux implémentations de la fonction sigmoïde :
  \begin{itemize}
    \item[A.] \texttt{return 1 / (1 + np.exp(-z))}
    \item[B.] \texttt{return np.where(z >= 0, 1 / (1 + np.exp(-z)), np.exp(z) / (1 + np.exp(z)))}
  \end{itemize}
}
\begin{enumerate}

\item
  \question{Les deux fonctions calculent-elles la même valeur mathématique ?}
  \reponse{
    Oui. Pour $z < 0$, on peut écrire :
    \[
      \frac{1}{1 + e^{-z}} = \frac{e^z}{e^z + 1}.
    \]
    Les deux expressions sont donc mathématiquement équivalentes.
  }

\item
  \question{Pour $z = -1000$, que se passe-t-il avec A ? Et avec B ?}
  \reponse{
    Avec A : \texttt{np.exp(1000)} produit un dépassement de capacité (\emph{overflow}) et retourne \texttt{inf},
    ce qui donne $\frac{1}{1 + \texttt{inf}} = 0$ — numériquement instable selon les implémentations.

    Avec B : pour $z = -1000 < 0$, on calcule \texttt{np.exp(-1000)} $\approx 0$,
    donc $\frac{e^{-1000}}{1 + e^{-1000}} \approx 0$, ce qui est numériquement stable.
  }

\item
  \question{Quel problème pratique B résout-il ?}
  \reponse{
    B évite l'overflow de \texttt{np.exp(-z)} pour $z \ll 0$ en utilisant une forme équivalente
    qui ne fait jamais appel à \texttt{exp} d'un grand nombre positif.
    C'est la technique classique du \emph{log-sum-exp} appliquée à la sigmoïde.
  }

\end{enumerate}
}
